%%%%%%%%%%%%%%%%%%%%%%%%%%%%%%%%%%%%%%%%%%%%%%%%%%%%%%%%%%%%%%%
% Based on a template by Dana Ernst.
% --------------------------------------------------------------

\documentclass[12pt]{article}

\usepackage[margin=1in]{geometry} 
\usepackage{amsmath,amsthm,amssymb}
%\usepackage{cancel}
\usepackage[makeroom]{cancel}
\usepackage{mathtools}
\usepackage[thinc]{esdiff}
\usepackage{tikz}
\usetikzlibrary{matrix}
\usepackage{amssymb}

\newcommand{\N}{\mathbb{N}}
\newcommand{\Z}{\mathbb{Z}}

\newenvironment{given}[1][Given]{\begin{trivlist}
\item[\hskip \labelsep {\bfseries #1:}\hskip \labelsep ]}{\end{trivlist}}

\newenvironment{derivation}[1][Derivation]{\begin{trivlist}
\item[\hskip \labelsep {\bfseries #1:}]}{\end{trivlist}}

\newcommand{\boxalign}[2][0.97\textwidth]{
 \par\noindent\tikzstyle{mybox} = [draw=black,inner sep=6pt]
 \begin{center}\begin{tikzpicture}
  \node [mybox] (box){%
   \begin{minipage}{#1}{\vspace{-5mm}#2}\end{minipage}
  };
 \end{tikzpicture}\end{center}
}


\begin{document}

\hfill {\Large \center Cross Correlaton of Periodic Signals Constructed from DFT Bases} \\
\\
\noindent Keith Rausch \\
\hfill \today \\

\newcommand\twopi{2 \pi}
\newcommand\twopii{\twopi \mathrm{i}}
\newcommand\pii{\pi \mathrm{i}}

\begin{given}{} 
Definition of Cross Correlation for complex signals:
\begin{align}
(f \star g)(\tau) & = \int_{-\infty}^{+\infty} \overline{f(t)} g(t+\tau) \,dt
\end{align}

For periodic signals with period of $T$:
\begin{align}
(f \star g)(\tau) & = \int_{t_0}^{t_0+T} \overline{f(t)} g(t+\tau) \,dt
\end{align}

Summation Properties:
\begin{align}
\left( \sum_{i} C_i \right) \left( \sum_{j} C_j \right) & = \sum_{i} \sum_{j} C_i C_j
\end{align}

Complex Conjugate Properties:
\begin{align}
\overline{zw} & = (\overline{z}) (\overline{v}) \\
\overline{ \exp{z}} & = \exp{\overline{z}}
\end{align}

%\begin{given}
Given 2 signals with $N$ samples, where $N$ is even:
\begin{align}
y_1(t) & = \sum_{m = -N/2}^{N/2-1} A_m \exp{[\twopii f_m t]} \\
y_2(t) & = \sum_{m = -N/2}^{N/2-1} B_m \exp{[\twopii f_m t]}
\end{align}
Where $d$ is the sampling period (in seconds per sample), and $N$ is the number of samples in a period $T$, we have $f_m = m/(dN) = m/T$
\end{given}

%\end{ex}

%\begin{sol}\

\begin{derivation}

Starting with the definition of convolution, we can substitute the $y_1$ and $y_2$ for their Fourier representations, and the integration bounds can be replaced with the period ($T$) of the given signals:
\begin{align}
(f \star g)(\tau) & = \int_{-\infty}^{+\infty} \overline{y_1(t)} y_2(t+\tau) \,dt \\
& = \int_{t_0}^{t_0+T} \overline{y_1(t)} y_2(t+\tau) \,dt \\
& =  \int_{t_0}^{t_0+T} \left( \sum_{m = -N/2}^{N/2-1} \overline{A_m \exp{[\twopii f_m t]}} \right) \left( \sum_{n = -N/2}^{N/2-1} B_n \exp{[\twopii f_n (t+\tau)]} \right) \,dt
\end{align}

Rearranging terms and moving the complex conjugate inside the exponential. The argument to $\exp{[\twopii f_m t]}$ is purely imaginary, so we can write $\overline{ \exp{[\twopii f_m t]} } = \exp{[\overline{ \twopii f_m t }]} = \exp{[-\twopii f_m t]}$:
\begin{align}
& =  \int_{t_0}^{t_0+T} \sum_{m = -N/2}^{N/2-1} \sum_{n = -N/2}^{N/2-1} \overline{A_m} \exp{[-\twopii f_m t]} B_n \exp{[\twopii f_n (t+\tau)]} \,dt \\
& =  \int_{t_0}^{t_0+T} \sum_{m = -N/2}^{N/2-1} \sum_{n = -N/2}^{N/2-1} \overline{A_m} B_n \exp{[\twopii f_n \tau]} \exp{[\twopii (f_n-f_m) t]} \,dt
\end{align}

The expression $f_n - f_m$ is $0$ where $n=m$, so we can't immediately integrate this expression. Instead, we must break the sums apart to isolate the division by $0$:
\begin{align}
& \!\begin{aligned}
 =  \int_{t_0}^{t_0+T} & \sum_{m = -N/2}^{N/2-1} \sum_{\substack{n = -N/2 \\ n \neq m}}^{N/2-1} \overline{A_m} B_n \exp{[\twopii f_n \tau]} \exp{[\twopii (f_n-f_m) t]} \\
 + & \sum_{m = -N/2}^{N/2-1} \sum_{\substack{n = -N/2 \\ n=m}}^{N/2-1} \overline{A_m} B_n \exp{[\twopii f_n \tau]} \exp{[\twopii \cancelto{0}{(f_n-f_m)} t ]} \,dt
\end{aligned}
\end{align}

We can move the integral inside since the summation bounds does not depend on $t$ (FACT CHECK THIS:
\begin{align}
& \!\begin{aligned}
 = & \sum_{m = -N/2}^{N/2-1} \sum_{\substack{n = -N/2 \\ n \neq m}}^{N/2-1} \overline{A_m} B_n \exp{[\twopii f_n \tau]} \int_{t_0}^{t_0+T} \exp{[\twopii (f_n-f_m) t]} \,dt \\
 + & \sum_{m = -N/2}^{N/2-1} \sum_{\substack{n = -N/2 \\ n=m}}^{N/2-1} \overline{A_m} B_n \exp{[\twopii f_n \tau]} \int_{t_0}^{t_0+T} 1 \,dt
\end{aligned}
\end{align}

Evaluating the integral:
\begin{align}
& \!\begin{aligned}
 = & \sum_{m = -N/2}^{N/2-1} \sum_{\substack{n = -N/2 \\ n \neq m}}^{N/2-1} \overline{A_m} B_n \exp{[\twopii f_n \tau]} \left( \frac{1}{ \twopii (f_n-f_m)} \right) \left[ \exp{[\twopii (f_n-f_m) t]} \right] \Big|_{t_0}^{t_0+T} \\
 + T & \sum_{m = -N/2}^{N/2-1} \sum_{\substack{n = -N/2 \\ n=m}}^{N/2-1} \overline{A_m} B_n \exp{[\twopii f_n \tau]}
\end{aligned}
\end{align}

\pagebreak

Since $t_0$ can be arbitrarily chosen, we set $t_0=0$. Since $i$ and $j$ are integers, $ \twopii (i + j) $ is simply a multiple of $2 \pi$, and the integral goes to zero:
\begin{align}
& \!\begin{aligned}
 = & \sum_{m = -N/2}^{N/2-1} \sum_{\substack{n = -N/2 \\ n \neq m}}^{N/2-1} \overline{A_m} B_n \exp{[\twopii f_n \tau]} \left( \frac{1}{ \twopii (f_n-f_m)} \right) \cancelto{0}{\left[ \exp{[\twopii \frac{i+j}{T} T]} - \exp{[\twopii \frac{i+j}{T} 0]} \right]} \\
 + T & \sum_{m = -N/2}^{N/2-1} \sum_{\substack{n = -N/2 \\ n=m}}^{N/2-1} \overline{A_m} B_n \exp{[\twopii f_n \tau]}
\end{aligned}
\end{align}

Since $n=m$, the summation can be simplified:
\begin{align}
& = T \sum_{m = -N/2}^{N/2-1} \sum_{\substack{n = -N/2 \\ n=m}}^{N/2-1} \overline{A_m} B_n \exp{[\twopii f_n \tau]} \\
& = T \sum_{m = -N/2}^{N/2-1} \overline{A_m} B_{m} \exp{[\twopii f_m \tau]}
\end{align}

Finally:
\boxalign[0.7\textwidth]{
\begin{align}
(y_1 \star y_2) (\tau) & = T \sum_{m = -N/2}^{N/2-1} \overline{A_m} B_{m} \exp{[\twopii f_m \tau]}
\end{align}
}

Differentiating with respect to $\tau$:
\boxalign[0.7\textwidth]{
\begin{align}
\diff{^{(n)}}{\tau^{(n)}} (y_1 \star y_2)(\tau) & = (\twopii)^{n} T \sum_{m = -N/2}^{N/2-1} f_m^n \overline{A_m} B_m \exp{[\twopii f_m \tau]}
\end{align}
}

Let's apply the ideas of Cooley and Tukey and decimate the FFT into even and odd components that can be reused. Remember that $f_m = \frac{m}{dN} = \frac{m}{T}$. For brevity, lets set $C_m = f_m^n \overline{A_m} B_m $
\begin{align}
\diff{^{(n)}}{\tau^{(n)}} (y_1 \star y_2)(\tau) \left(\frac{1}{(\twopii)^{n} T}\right) & = \sum_{m = -N/4}^{N/4} C_{2m} \exp{[\twopii f_{2m} \tau]} \\
& + \sum_{m = -N/4}^{N/4-1} C_{2m+1} \exp{[\twopii f_{2m+1} \tau]} \\
\end{align}

Since $f_m=m/T$, then $\exp{[\twopii f_{2m+1} \tau]} = \exp{[\twopii \frac{2m+1}{T} \tau]} = \exp{[\twopii \frac{2m}{T} \tau]} \exp{[\twopii \frac{\tau}{T}]} $ and you can factor out $\exp{[\twopii \frac{\tau}{T}]}$ from the summation:
\begin{align}
\diff{^{(n)}}{\tau^{(n)}} (y_1 \star y_2)(\tau) \left(\frac{1}{(\twopii)^{n} T}\right) = & \sum_{m = -N/4}^{N/4} C_{2m} \exp{[\twopii f_{2m} \tau]} \\
+ \exp{[\twopii \frac{\tau}{T}]} & \sum_{m = -N/4}^{N/4-1} C_{2m+1} \exp{[\twopii f_{2m} \tau]} \\
= & \sum_{m = -N/4}^{N/4} C_{2m} \exp{[\twopii 2 f_m \tau]} \\
+ \exp{[\twopii \frac{\tau}{T}]} & \sum_{m = -N/4}^{N/4-1} C_{2m+1} \exp{[\twopii 2 f_m \tau]}
\end{align}

Let's set $E(\tau)$ and $O(\tau)$ to hold the even and odd terms respectively
\begin{align}
E(\tau) & = \sum_{m = -N/4}^{N/4} C_{2m} \exp{[\twopii 2 f_m \tau]} \\
O(\tau) & = \exp{[\twopii \frac{\alpha}{T}]} \sum_{m = -N/4}^{N/4-1} C_{2m+1} \exp{[\twopii 2 f_m \tau]}
\end{align}


Now Let's check that  $\diff{^{(n)}}{\tau^{(n)}} (y_1 \star y_2)(\tau) = \diff{^{(n)}}{\tau^{(n)}} (y_1 \star y_2)(\tau+\alpha)$. Remember that $T = dN$.
\begin{align}
\diff{^{(n)}}{\tau^{(n)}} (y_1 \star y_2)(\tau+\alpha) \left(\frac{1}{(\twopii)^{n} T}\right)
= & \sum_{m = -N/4}^{N/4} C_{2m} \exp{[\twopii 2 f_m \tau]} \exp{[\twopii 2 f_m \alpha]} \\
+ \exp{[\twopii \frac{\tau}{T}]} \exp{[\twopii \frac{\alpha}{T}]} & \sum_{m = -N/4}^{N/4-1} C_{2m+1} \exp{[\twopii 2 f_m \tau]} \exp{[\twopii 2 f_m \alpha]}
\end{align}

Let's set $\alpha = T/2$, so $2 f_m \alpha = \frac{2m}{dN} \frac{dN}{2} = m$ and $\frac{\alpha}{T} = \frac{1}{2}$.
\begin{align}
\diff{^{(n)}}{\tau^{(n)}} (y_1 \star y_2)(\tau+\alpha) \left(\frac{1}{(\twopii)^{n} T}\right) = & \sum_{m = -N/4}^{N/4} C_{2m} \exp{[\twopii 2 f_m \tau]} \exp{[\twopii 2 f_m \alpha]} \\
+ \exp{[\twopii \frac{\tau}{T}\tau]} \exp{[\twopii \frac{\alpha}{T}]} & \sum_{m = -N/4}^{N/4-1} C_{2m+1} \exp{[\twopii 2 f_m \tau]} \exp{[\twopii 2 f_m \alpha]} \\
= & \sum_{m = -N/4}^{N/4} C_{2m} \exp{[\twopii 2 f_m \tau]} \cancelto{1}{\exp{[\twopii m]}} \\
+ \exp{[\twopii \frac{\tau}{T}]} \cancelto{-1}{\exp{[\pii]}} & \sum_{m = -N/4}^{N/4-1} C_{2m+1} \exp{[\twopii 2 f_m \tau]} \cancelto{1}{\exp{[\twopii 2 f_m \alpha]}} \\
= & \sum_{m = -N/4}^{N/4} C_{2m} \exp{[\twopii 2 f_m \tau]} \\
- \exp{[\twopii \frac{\tau}{T}]} & \sum_{m = -N/4}^{N/4-1} C_{2m+1} \exp{[\twopii 2 f_m \tau]}
\end{align}

Finally, the butterfly appears:
\begin{align}
\diff{^{(n)}}{\tau^{(n)}} (y_1 \star y_2)(\tau) \left(\frac{1}{(\twopii)^{n} T}\right) & = E(\tau) + \exp{[\twopii \frac{\tau}{T}]} O(\tau) \\
\diff{^{(n)}}{\tau^{(n)}} (y_1 \star y_2)(\tau+\frac{T}{2}) \left(\frac{1}{(\twopii)^{n} T}\right) & = E(\tau) - \exp{[\twopii \frac{\tau}{T}]} O(\tau)
\end{align}


%%%%%%%%%%%%%%%%%%%%%%%%%%%%%%%%%%%%%%%%%%%%%%%%%%%%%%%%%%%%%%%%%%%%%%%
\newpage
You can arrive at the same conclusion by converting to the frequency domain:
\begin{align}
(f \star g)(\tau) & = \int_{-\infty}^{+\infty} \overline{y_1(t)} y_2(t+\tau) \,dt \\
\mathfrak{F}[f \star g](\tau) & =  \overline{\mathfrak{F}\left[y_1(t)\right]} \mathfrak{F}\left[y_2(t+\tau)\right] \\
\mathfrak{F}[f \star g](\tau) & =  \overline{\mathfrak{F}\left[\left( \sum_{m = -N/2}^{N/2-1} A_m \exp{[\twopii f_m t]} \right)\right]} \mathfrak{F}\left[\left( \sum_{n = -N/2}^{N/2-1} B_n \exp{[\twopii f_n (t+\tau)]} \right)\right] \\
\mathfrak{F}[f \star g](\tau) & =  \left( \sum_{m = -N/2}^{N/2-1} \overline{A_m} \twopi \delta (\omega - \twopi f_m) \right) \left( \sum_{n = -N/2}^{N/2-1} B_n \twopi  \delta(\twopii f_n) \exp{[-\twopii f_n \tau]} \right) \\
\mathfrak{F}[f \star g](\tau) & =  \sum_{m = -N/2}^{N/2-1} \sum_{n = -N/2}^{N/2-1} \overline{A_m} B_n \twopi  \delta (\omega - \twopi f_m) \twopi  \delta \left(\omega - \twopii f_n \right) \exp{[-\twopii f_n \tau]} \\
\mathfrak{F}[f \star g](\tau) & =  \sum_{m = -N/2}^{N/2-1} \overline{A_m} B_m \twopi \delta (\omega - \twopi f_m)  \twopi \delta \left(\omega - \twopii f_m \right) \exp{[-\twopii f_n \tau]} \\
\mathfrak{F}[f \star g](\tau) & = (\twopi )^2 \sum_{m = -N/2}^{N/2-1} \overline{A_m} B_m \delta (\omega - \twopi f_m)  \exp{[-\twopii f_n \tau]}
\end{align}


\end{derivation}
%\end{sol}

% --------------------------------------------------------------
%     You don't have to mess with anything below this line.
% --------------------------------------------------------------

\end{document}





%XXXX:
%\begin{figure}[!h]
%\centering
%\begin{tikzpicture} [
%    toplabel/.style={
%        row 1 column #1/.style={%
%            execute at begin cell={%
%                |[label=\the\pgfmatrixcurrentcolumn]|},
%            execute at empty cell={%
%                \node[label=\the\pgfmatrixcurrentcolumn]{};}
%        }
%    },
%    leftlabel/.style={
%        row #1 column 1/.style={%
%            execute at begin cell={%
%                |[label=left:\the\pgfmatrixcurrentrow]|},
%            execute at empty cell={%
%                \node[label=left:\the\pgfmatrixcurrentrow]{};}
%        }
%    },
%    topleft/.style={
%        row 1 column 1/.style={%
%            execute at begin cell={%
%                |[label=\the\pgfmatrixcurrentcolumn,
%                    label=left:\the\pgfmatrixcurrentrow]|},
%            execute at empty cell={%
%                \node[label=left:\the\pgfmatrixcurrentrow,
%                    label=\the\pgfmatrixcurrentcolumn]{};}
%        }
%    }
%]
%\matrix (A) [
%    matrix of nodes, 
%    nodes in empty cells,
%    nodes={draw, minimum size=1cm, anchor=center},
%    row sep=-\pgflinewidth,
%    column sep=-\pgflinewidth,
%    topleft,
%    toplabel/.list={2,3,4,5,6,7,8},
%    leftlabel/.list={2,...,8}
%]
%{x& & & & & & & \\
% & & & & & & & \\
% & & & & & & & \\
% & & & & & & & \\
% & & & & & & & \\
% & & & & & & & \\
% & & & & & & & \\
% & & & & & & & \\};
%
%\end{tikzpicture}
%%\tikz
%\caption{A caption.}
%\end{figure}
